% Standalone Overleaf document. Compiler: pdfLaTeX.
% No external figures, bibliography, source files, or shell escape are required.
\documentclass[11pt,a4paper]{article}
\usepackage[T1]{fontenc}
\usepackage[utf8]{inputenc}
\usepackage{lmodern}
\usepackage[margin=25mm]{geometry}
\usepackage{amsmath,amssymb,booktabs,longtable,tabularx}
\usepackage{xcolor,listings,hyperref,microtype}
\hypersetup{colorlinks=true,linkcolor=blue!55!black,urlcolor=blue!55!black,
  pdftitle={SpectralMM: A User Guide for the Python and R Interfaces}}
\definecolor{codebackground}{RGB}{247,248,250}
\lstset{basicstyle=\ttfamily\footnotesize,backgroundcolor=\color{codebackground},
  frame=single,rulecolor=\color{black!15},breaklines=true,columns=fullflexible,
  keepspaces=true,showstringspaces=false,aboveskip=8pt,belowskip=8pt}
\setlength{\emergencystretch}{2em}
\setlength{\parindent}{0pt}
\setlength{\parskip}{5pt}
\newcommand{\pkg}{\textsf{SpectralMM}}
\newcommand{\code}[1]{\texttt{#1}}
\title{\pkg: A User Guide for the Python and R Interfaces}
\author{}
\date{Version 0.1.0}
\begin{document}
% R installation is independent of the standalone CMake build.

\maketitle
\begin{abstract}
\pkg{} fits generalized linear models and residual-based statistical models
using a shared C++ computational core. It provides spectral-preconditioned
PCG and MM solvers, direct Cholesky, and five alternative Krylov solvers for dense and sparse design
matrices. This guide describes installation from source, the Python and R
interfaces, supported models, convergence diagnostics, and reproducible
performance comparisons with the original Julia implementation.
\end{abstract}
\tableofcontents
\clearpage

\section*{Companion small-data examples}
The standalone document \texttt{Small\_Model\_Examples.tex} provides runnable
100-by-3 examples for all fourteen models, fitted coefficients and inference,
three-language comparisons, and complete Python iteration logs. The source
scripts and results are in the repository's \texttt{examples} directory.

\section*{Automatic Python and Julia installation}
From a local checkout, run \texttt{python -m pip install .}; subsequent Python
sessions use \texttt{import spectralmm}. The wheel contains the shared library.
In Julia, use \texttt{Pkg.develop(path="/path/to/SpectralMM")} and then
\texttt{using SpectralMM}. Call \texttt{SpectralMM.fit(X, y; family=:bernoulli, backend=:cpp)}
for the C++ implementation. Omitting \texttt{backend} selects pure Julia.
Both backends support GLM and non-GLM models through the same matrix API.
The returned model exposes \texttt{model.backend}. On first C++ use, the package builds the core automatically and reuses
its cached library across sessions. All fourteen families are supported.
Existing positional-family Julia methods and \texttt{glm} retain their original
implementation. The manual CMake instructions below are for standalone development.
See \texttt{docs/Automatic\_Installation.md} for details.

\section*{Automatic R package installation}
The repository root is a standard R package. After the updated source is
published, install it with \texttt{remotes::install\_github("Xunjian-Li/SpectralMM")}.
For a local checkout use \texttt{remotes::install\_local("/path/to/SpectralMM")}.
Installation compiles the shared C++ core automatically using R's build system,
RcppEigen and R's configured BLAS. Subsequent sessions only need
\texttt{library(SpectralMM)}. A C++17 toolchain is required for source installation;
R users do not need the standalone CMake commands below. Formula and matrix interfaces are now available in all three languages; see
\texttt{docs/Statistical\_API.md} for the preferred statistical API. See \texttt{docs/R\_Package.md} for the complete workflow.

\section{Overview}
The Python and R interfaces call the same C++ library. Optimization iterations
run inside C++; neither interface requires a Julia runtime for model fitting.
Julia is needed only when running the original Julia implementation or generating
the shared reference data used in the supplied benchmarks.

The repository provides a standard R package, Python source and wheel builds,
and a Julia package with automatic C++ build/cache support. The C++ interfaces
support fourteen families. Original Julia algorithms remain available.
GitHub installation requires publishing these changes; no public index release
is claimed. Repository scripts below assume its root as the working directory.

\subsection{Computational approach}
For a feature matrix $X\in\mathbb{R}^{n\times d}$, slope vector
$\beta\in\mathbb{R}^{d}$, and intercept $\alpha$, the linear predictor is
$\eta=\alpha\mathbf{1}+X\beta$. The total parameter count is $p=d+1$
with an intercept and $p=d$ without one.
The solver minimizes a sum of model-specific losses, optionally with a ridge penalty:
\[
 F(\alpha,\beta)=\sum_{i=1}^{n}\ell_i(\alpha+x_i^\top\beta;y_i)
       +\frac{\lambda}{2}\|\beta\|^2.
\]
The automatic intercept is excluded from ridge by default. Set
\code{penalize\_intercept} to true to penalize it as well. Response-only
likelihood constants may be omitted from the reported loss.

Weighted least-squares subproblems use an operator of the form
\[
 H_Wv=\widetilde X^\top\!\left(W(\widetilde Xv)\right)+\lambda Dv,
 \qquad \widetilde X=[\mathbf{1},X],\quad D=\operatorname{diag}(0,1,\ldots,1).
\]
The intercept column is implemented implicitly in the C++ products, without
copying the feature matrix. This operator is evaluated without forming a dense
$p\times p$ Hessian. Without an intercept, use $\widetilde X=X$ and $D=I$.
Depending on the family, $H_W$ represents a Hessian, Fisher-scoring operator, or
MM surrogate curvature. In particular, Student-$t$ uses the positive MM weights
of the Julia implementation rather than the exact Hessian of its loss.

The spectral solver uses a small Lanczos subspace and spectral correction to
construct its preconditioner. Cached products such as $HQ$ are retained without
materializing the full Hessian. The parameter \code{floor} controls the spectral
regularization parameter $\rho$; it is distinct from the ridge penalty $\lambda$.

\subsection{Supported solvers}
\begin{center}
\begin{tabular}{ll}
\toprule
\code{solver} & Method \\
\midrule
\code{spectral}/\code{pcg} & Spectral-preconditioned conjugate gradient (PCG) \\
\code{mm} & Spectral MM with quadratic line search \\
\code{cho}/\code{cholesky} & Direct Cholesky weighted least squares \\
\code{cg} & Conjugate gradient \\
\code{cgls} & Conjugate gradient for least squares \\
\code{crls} & Conjugate residual for least squares \\
\code{lsqr} & LSQR \\
\code{lsmr} & LSMR \\
\bottomrule
\end{tabular}
\end{center}
All eight methods support the same native model families and dense/CSC inputs.
Cholesky forms a full curvature matrix; the other methods use matrix-vector products.
Original Julia high-level fitting supports the same choices as symbols, such as
\code{solver=:cho} or \code{solver=:mm}. The Julia C++ interface supports all fourteen families. The low-level \code{spectral\_mm} function
supports only \code{:pcg} and \code{:mm}. CG uses Jacobi scaling by default;
set \code{preconditioner="none"} (Julia \code{:none}) for unpreconditioned CG.
Cholesky retries failed factorizations with small diagonal jitter. This numerical
safeguard does not alter the objective or remove inference validity checks.
Archived native benchmark tables still cover the original six choices.

\section{Installation}
\subsection{Build the shared library}
The native build requires a C++17 compiler, CMake 3.18 or later, and Eigen 3.4
headers. Replace the example Eigen path with a directory containing
\path{Eigen/Core}.
\begin{lstlisting}[language=bash]
export EIGEN3_INCLUDE_DIR=/path/to/eigen-3.4.0
cmake -S cpp -B build/native \
  -DCMAKE_BUILD_TYPE=Release \
  -DEIGEN3_INCLUDE_DIR="$EIGEN3_INCLUDE_DIR" \
  -DSPECTRALMM_USE_BLAS=ON
cmake --build build/native -j2
\end{lstlisting}
This command enables system BLAS for supported dense Eigen operations.
The current macOS binary links Apple Accelerate; add
\texttt{-DBLA\_VENDOR=Apple} to select it explicitly.
Restart language sessions after replacing a loaded library.
Rebuild the library after changing the native source or moving to another machine.

\subsection{Python dependencies}
NumPy is required for fitting. SciPy is required for statistical inference, sparse input and the examples
below. The full benchmark additionally uses statsmodels and threadpoolctl.
\begin{lstlisting}[language=bash]
python3 -m pip install .
# Optional benchmark dependencies:
python3 -m pip install statsmodels threadpoolctl
\end{lstlisting}

\subsection{R dependencies}
Install the standard R package once; Rcpp and RcppEigen are resolved automatically.
Sparse examples use Matrix. The standalone benchmark also uses jsonlite and MASS.
\begin{lstlisting}[language=R]
install.packages("remotes")
remotes::install_github("Xunjian-Li/SpectralMM")
library(SpectralMM)
\end{lstlisting}
The installation compiles and registers the native code. No manual loading is
needed during a session beyond \code{library(SpectralMM)}.

\subsection{Library discovery}
The installed Python package loads its bundled library; Julia automatically
builds and reuses a library in its scratch cache. Neither requires a manual path.
The installed R package instead loads its own registered library automatically.
The expected names are \code{libspectralmm.dylib} on macOS,
\code{libspectralmm.so} on Linux, and \code{spectralmm.dll} on Windows.
To use another build, set \code{SPECTRALMM\_LIBRARY} to its absolute path before
loading the interface. The shell examples in this guide use POSIX syntax.
The retained validation results were obtained on macOS arm64; they do not
constitute testing on every operating system.

\section{Quick Start}
\subsection{Python: logistic regression}
The interface fits an intercept automatically. Supply feature columns only;
To fit without an automatic intercept, set the option below. It also supports
complete design matrices that already contain a constant column.
\begin{lstlisting}
fit_intercept=False
\end{lstlisting}
\begin{lstlisting}[language=Python]
import numpy as np
from scipy.special import expit
from spectralmm import fit

rng = np.random.default_rng(3124)
Z = rng.normal(size=(200, 6))
X = Z
beta = np.array([0.5, 0.2, -0.3, 0.1, 0.2, -0.1, 0.15])
y = (rng.random(200) < expit(beta[0] + X @ beta[1:])).astype(float)

model = fit(X, y, family="bernoulli", solver="spectral",
            rank=5, floor=1e-6)
print(model.intercept_)
print(model.coef_)
print(model.info)
probability = model.predict(X)
\end{lstlisting}

\subsection{R: logistic regression}
\begin{lstlisting}[language=R]
library(SpectralMM)

set.seed(3124)
Z <- matrix(rnorm(200 * 6), nrow=200)
X <- Z
beta <- c(0.5, 0.2, -0.3, 0.1, 0.2, -0.1, 0.15)
y <- rbinom(200, 1, plogis(as.vector(beta[1] + X %*% beta[-1])))

model <- spectralmm_fit(X, y, family="bernoulli",
                        solver="spectral", rank=5, floor=1e-6)
coef(model)
model$info
probability <- predict(model, X)
\end{lstlisting}
The Python and R examples are independent simulations. Identical seed values do
not make their random samples identical. Use the shared binary benchmark inputs
when comparing the interfaces.

\subsection{Intercept conventions and Julia usage}
Python uses \code{fit\_intercept=True}; R and Julia use
\code{intercept=TRUE} and \code{intercept=true}, respectively.
Prediction accepts the same feature columns as training and adds the fitted
intercept automatically. R labels the coefficient \code{(Intercept)}.
Do not add a column of ones when automatic intercept is enabled. A nonzero
constant feature triggers a warning; it is not silently removed.

To retain an existing explicit design matrix, disable automatic intercept.
All supplied columns then retain the legacy ridge penalty. This is the convention
used by the benchmark scripts, whose fixtures already contain a constant column.
The low-level C ABI and original Julia \code{spectral\_mm} continue to accept
complete design matrices. The original Julia low-level solvers can exclude the
first column from ridge using \code{penalize\_intercept=false}; for
\code{spectral\_mm}, set this flag in \code{SpectralOptions}.

The original Julia package now uses the same matrix-interface convention:
\begin{lstlisting}
using SpectralMM, Distributions, GLM
model = SpectralMM.glm(X, y, Bernoulli(), LogitLink();
                       intercept=true, ridge=0.1)
model = SpectralMM.fit(X, y, PseudoHuber(1.0); intercept=true)
\end{lstlisting}
Its formula interface controls the intercept through the formula:
\code{y ~ x1 + x2} includes it, \code{y ~ 0 + x1 + x2} excludes it,
and \code{y ~ 1} fits only an intercept. Do not also pass the
\code{intercept} keyword to formula fits. The standalone native Julia wrapper
uses \code{SpectralMM.fit(X, y; backend=:cpp, intercept=true)}.
Python, R, and both Julia backends now support formula input with stored categorical encoding.

\subsection{Residual models}
Model-specific parameters are supplied through \code{family\_options}, using a
Python dictionary or a named R list.
\begin{lstlisting}[language=Python]
y_cont = beta[0] + X @ beta[1:] + rng.standard_t(4, size=X.shape[0])
robust = fit(X, y_cont, family="student_t",
             family_options={"nu": 4, "sigma": 1})
quantile = fit(X, y_cont, family="smooth_quantile",
               family_options={"tau": 0.25, "smoothing": 0.1})
\end{lstlisting}
\begin{lstlisting}[language=R]
y_cont <- as.vector(beta[1] + X %*% beta[-1]) + rt(nrow(X), df=4)
robust <- spectralmm_fit(X, y_cont, family="student_t",
                        family_options=list(nu=4, sigma=1))
quantile <- spectralmm_fit(
  X, y_cont, family="smooth_quantile",
  family_options=list(tau=0.25, smoothing=0.1))
\end{lstlisting}

\subsection{Sparse matrices and alternative solvers}
Python accepts SciPy sparse matrices and converts them to canonical CSC format
when necessary. R accepts \code{dgCMatrix} objects.
\begin{lstlisting}[language=Python]
from scipy import sparse
Xs = sparse.csc_matrix(X)
model_cg = fit(Xs, y, family="bernoulli", solver="cg",
               preconditioner="jacobi")
\end{lstlisting}
\begin{lstlisting}[language=R]
library(Matrix)
Xs <- as(Matrix(X, sparse=TRUE), "dgCMatrix")
model_cg <- spectralmm_fit(Xs, y, family="bernoulli",
                          solver="cg", preconditioner="jacobi")
\end{lstlisting}
Changing a dense matrix to CSC does not create sparsity; this example illustrates
the accepted input format. Sparse computations are most useful when the design
already contains relatively few nonzero entries.

\section{Model Families}
Both interfaces use the names below. Unknown or inapplicable model parameters
are rejected. All responses must be finite.
\begin{center}
\small
\begin{tabular}{lll}
\toprule
\code{family} & Model & Parameter defaults \\
\midrule
\code{gaussian} & Gaussian identity & None \\
\code{bernoulli} & Bernoulli logit & None \\
\code{probit} & Bernoulli probit & None \\
\code{poisson} & Poisson log & None \\
\code{gamma} & Gamma log & None \\
\code{negative\_binomial} & Negative-binomial log & \code{theta=1} \\
\code{smooth\_quantile} & Smooth quantile & \code{tau=0.5, smoothing=0.1} \\
\code{expectile} & Expectile & \code{tau=0.5} \\
\code{pseudo\_huber} & Pseudo-Huber & \code{delta=1.0} \\
\code{student\_t} & Student-$t$ & \code{nu=4, sigma=1} \\
\code{gaussian\_log} & Gaussian log & None \\
\code{gamma\_inverse} & Gamma inverse & None \\
\code{tweedie} & Tweedie log & \code{power=1.5} \\
\code{binomial} & Binomial logit & \code{trials=1} \\
\bottomrule
\end{tabular}
\end{center}
The \code{smoothing} parameter corresponds to Julia's \code{epsilon}.
Require $0<\tau<1$, positive finite scale/shape/trial-count parameters, and
$1\leq\code{power}\leq2$.

Bernoulli responses must be exactly zero or one. Poisson, negative-binomial,
and Tweedie responses must be nonnegative. Gamma, Gaussian-log, and Gamma-inverse
responses must be positive. Binomial responses must satisfy
$0\leq y_i\leq n_i$, where \code{trials} specifies a scalar or a length-$n$ vector.
The four residual models accept finite real responses.

For Gamma inverse, the initial linear predictor must be positive. Supply a
feasible \code{beta0} when automatic initialization cannot produce one, for
example when there is no constant column. Probit follows the normal-CDF
approximation used in the Julia implementation.

\section{Function Reference}
\subsection{Fitting a model}
\begin{lstlisting}[language=Python]
fit(X, y, family="gaussian", *, family_options=None,
    beta0=None, fit_intercept=True, penalize_intercept=False,
    solver="spectral", preconditioner="jacobi",
    krylov_rtol=1e-4, krylov_atol=1e-8, trace=False, verbose=False,
    **options)
\end{lstlisting}
\begin{lstlisting}[language=R]
spectralmm_fit(X, y, family="gaussian", beta0=NULL,
               intercept=TRUE, penalize_intercept=FALSE,
               solver="spectral", preconditioner="jacobi",
               krylov_rtol=1e-4, krylov_atol=1e-8,
               family_options=list(), ...)
\end{lstlisting}
These are abbreviated signatures; stopping options are described below.
Python uses keyword arguments for solver options; R passes them through
\code{...}. Python and R both support \code{verbose} and \code{trace}, disabled by default.
Use \code{verbose=True} in Python or \code{verbose=TRUE} in R to print the
iteration table. Enable \code{trace} to retain structured rows in the model.

\begin{description}
\item[\code{X}] A finite, real $n\times d$ feature matrix, with $n>0$.
Zero feature columns are allowed when fitting an intercept; otherwise $d\geq1$.
The core uses double precision. Conversion copies may be needed for dtype,
storage layout, or sparse normalization.
\item[\code{y}] A real vector of length $n$, satisfying the family-specific domain.
\item[\code{beta0}] An optional length-$p$ initial coefficient vector, with
the intercept first when enabled. Otherwise, the intercept is initialized from
the response mean and link; without one, coefficients start at zero when feasible.
\item[\code{family\_options}] Model parameters from the family table.
\item[\code{preconditioner}] Either \code{jacobi} or \code{none} for the five
alternative Krylov methods. The spectral solver uses its own preconditioner.
\end{description}

\subsection{Common numerical options}
\small
\begin{longtable}{@{}p{.24\linewidth}p{.16\linewidth}p{.54\linewidth}@{}}
\toprule
Option & Default & Meaning \\
\midrule
\endhead
\code{rank} & 0 & Automatic spectral rank: $p-1$ for $p<20$, otherwise 10.
An explicit rank must satisfy $1\leq k<p$. With $p=1$, automatic rank is zero. \\
\code{krylovdim} & 0 & Automatic initial Lanczos dimension. An explicit value
must lie in $[k+1,p]$. Restart attempts may increase the dimension. \\
\code{floor} & $10^{-6}$ & Spectral regularization parameter $\rho$. \\
\code{ridge} & 0 & Ridge penalty $\lambda$, excluding the automatic intercept by default. \\
\code{maxiter} & 200 & Maximum number of outer iterations. \\
\code{inner\_maxiter} & 100 & Maximum number of inner iterations. \\
\code{gtol} & $10^{-7}$ & Absolute gradient stopping tolerance. \\
\code{relgtol} & $10^{-8}$ & Relative gradient stopping tolerance. \\
\code{eta\_max} & 0.5 & Inner forcing parameter for spectral PCG. \\
\code{correction\_tol} & 0.01 & Threshold for projected curvature changes. \\
\code{resid\_tol} & 0.5 & Spectral residual acceptance tolerance. \\
\code{nesterov} & True & Enable outer Nesterov extrapolation. Use
\code{True}/\code{False} in Python and \code{TRUE}/\code{FALSE} in R. \\
\code{krylov\_rtol} & $10^{-4}$ & Relative inner tolerance for alternative
Krylov solvers. \\
\code{krylov\_atol} & $10^{-8}$ & Absolute inner tolerance for alternative
Krylov solvers. \\
\bottomrule
\end{longtable}
\normalsize
The benchmark settings differ from several interface defaults; they are listed
in Section~\ref{sec:benchmark}.

\subsection{Unified iteration table}
All solver choices use the columns Iter, Loss, RelGrad, Inner, InnerRes, EigRes,
Step and Spectrum. Row 0 describes the initial coefficients and has no inner
iteration or step. PCG/MM label the initial spectrum as \code{initial}; methods
without spectral information show a dash. Subsequent rows report loss and
relative gradient at the updated accepted coefficients. The final row and footer
use identical relative-gradient values and precision. Checking convergence at
the next loop entry is not an additional parameter update.

InnerRes is the relative WLS gradient residual before the outer line search;
it differs from the outer RelGrad. Spectrum reports correction, reuse, restart
or correction failure. CG, CHO and other non-spectral methods show dashes for
Spectrum and EigRes. The footer distinguishes gradient convergence from other
stopping reasons. Native output is buffered until the fit completes. Detailed
logging can add diagnostic matrix products; disable it for benchmarks.
Original Julia high-level functions use \code{verbose=true} for the same table.
The native Julia wrapper also offers \code{trace=true}; original Julia retains
its existing history arrays.

\subsection{Returned model and predictions}
Python returns a \code{Model} with \code{coef} (all parameters, intercept first),
\code{coef\_} (slopes only), \code{intercept\_} (scalar), \code{family},
\code{family\_options}, \code{info}, and an optional \code{trace}.
R returns an object of class \code{spectralmm\_native}, with \code{coef()},
\code{predict()}, and the corresponding list components.

Predictions are response means for GLMs and fitted locations, quantiles, or
expectiles for residual models. Binomial predictions are expected counts, not
success probabilities. For new binomial observations, pass appropriate trial
counts to \code{predict}:
\begin{lstlisting}[language=Python]
counts = binomial_model.predict(X_new, trials=10)
\end{lstlisting}
\begin{lstlisting}[language=R]
counts <- predict(binomial_model, X_new, trials=10)
\end{lstlisting}
The examples assume that \code{binomial\_model} was fitted with
\code{family="binomial"}. A vector of new trial counts must have one entry per
new observation.

\subsection{Diagnostics and stopping rules}
The \code{info} object reports outer and inner iteration counts, spectral restarts
and corrections, loss, absolute and relative gradient norms, spectral residual,
and termination status. In Python use \code{model.info["termination\_reason"]};
in R use \code{model\$info\$termination\_reason}.

With $g_0$ denoting the gradient at the initial coefficients, gradient convergence
means
\[
 \|g\|\leq\code{gtol}
 \quad\text{or}\quad
 \frac{\|g\|}{1+\|g_0\|}\leq\code{relgtol}.
\]
The wrappers also accept negligible steps and near-tolerance stalled steps by
default. The corresponding defaults are shown below; use R logical syntax
when calling the R interface. The step-relative tolerance is the square root
of double-precision machine epsilon.
\begin{lstlisting}
accept_negligible = True
accept_stalled = True
negligible_step_tol = 1e-14
step_reltol = 1.4901161193847656e-8
stalled_relgtol = 1e-4
\end{lstlisting}

Consequently, \code{converged} and \code{gradient\_converged} are distinct.
A fit accepted because of a small step need not satisfy the gradient tolerance.
Check both flags when comparing solvers at a common accuracy target.
The termination reasons are \code{gradient}, \code{maxiter},
\code{line\_search\_failed}, \code{inner\_breakdown},
\code{negligible\_step}, and \code{stalled\_step}.
Invalid inputs raise errors; ordinary nonconvergence is returned in the diagnostics.

\section{Optional Statistical Inference}
The interfaces support coefficient standard errors, full covariance matrices,
coefficient-wise Wald statistics, two-sided p-values, and confidence intervals.
Inference is a separate post-fit stage; it does not alter the fitted coefficients.

\subsection{Automatic size policy}
The default is \code{inference="auto"} in Python/R and
\code{inference=:auto} in Julia. Inference is attempted when the total parameter
count $p$ is at most \code{inference\_max\_p=50}. This count includes the
intercept and any expanded categorical or interaction columns. Above the cutoff,
the covariance allocation is skipped. Boolean true bypasses the size cutoff;
boolean false disables the post-fit stage entirely.

The cutoff is an engineering default for this implementation, not a universal
statistical rule or a common documented default of R \code{glm}, statsmodels,
and GLM.jl. One double-precision covariance matrix needs $8p^2$ bytes: about
20 kB at $p=50$ and 200 MB at $p=5000$. Factorization and sandwich calculations
need additional matrices. Dense work is approximately $O(np^2+p^3)$;
observation count and sparsity also affect runtime. The fitting algorithm remains
matrix-free, but this optional stage explicitly forms full statistical matrices.
It does not interpret the spectral preconditioner as an approximate covariance.

\subsection{Accessors}
\begin{lstlisting}[language=Python]
model = fit(X, y, family="bernoulli", inference="auto",
            inference_max_p=50, level=.95)
print(model.summary())
model.bse
model.pvalues
model.cov_params()
model.conf_int()
model.inference["wald_chisq"]
model.inference["wald_p_value"]
\end{lstlisting}
\begin{lstlisting}[language=R]
model <- spectralmm_fit(X, y, family="bernoulli",
                        inference="auto", inference_max_p=50)
summary(model)
vcov(model)
confint(model, level=.95)
model$inference$wald_chisq
\end{lstlisting}
\begin{lstlisting}
using SpectralMM, Distributions, GLM, StatsAPI
model = SpectralMM.glm(X, y, Bernoulli(), LogitLink();
                       inference=:auto, inference_max_p=50)
stderror(model)
vcov(model)
confint(model; level=.95)
coeftable(model)
model.inference.wald_chisq
\end{lstlisting}
The native Julia wrapper exposes its own \code{stderror}, \code{vcov}, and
\code{confint} functions and now uses Distributions.jl for reference distributions.
Python needs SciPy when inference is computed. Original Julia's coefficient table
is a named tuple of column vectors. Summary displays show at most 20 rows by
default; this display limit does not limit the stored inference results.

\subsection{Covariance and reference distributions}
Let $A$ denote the complete design, including the intercept. For GLMs, the default
model-based covariance is
\[
 \widehat{\operatorname{Var}}(\widehat\beta)
 =\widehat\phi(A^\top W A)^{-1},\qquad
 W_{ii}=\frac{(d\mu_i/d\eta_i)^2}{V(\mu_i)}.
\]
Inference uses statistical weights without the optimization weight floor.
Dispersion is fixed at one for Bernoulli, Probit, Binomial, Poisson, and
negative binomial with supplied shape. Gaussian, Gamma, and Tweedie models use
the Pearson dispersion estimate divided by $n-p$ unless a positive
\code{dispersion} value is supplied. Binomial trial counts enter the weights;
$n$ counts independent rows.

For residual models, \code{cov\_type="auto"} selects an HC1 sandwich estimator:
\[
 \widehat{\operatorname{Var}}(\widehat\beta)
 =\frac{n}{n-p}H^{-1}BH^{-1},\quad
 H=\sum_i x_i x_i^\top\frac{\partial q_i}{\partial\eta_i},\quad
 B=\sum_i x_i x_i^\top q_i^2.
\]
Here $q_i$ is the score of the fitted objective. Student-$t$ uses the actual score
derivative rather than its positive MM weights. SmoothQuantile inference targets
the smoothed objective at its fixed smoothing parameter. Family shape, scale, and
tuning parameters are treated as fixed. Sandwich covariance can also be requested
for GLMs; its bread uses observed score derivatives. No cluster or time-series
correction is included. Model-based covariance is not offered for residual models,
and a dispersion override is not applicable to sandwich covariance.

The coefficient statistic is $\widehat\beta_j/\operatorname{SE}_j$. Estimated
dispersion with model-based covariance defaults to a Student-$t_{n-p}$ reference;
other cases default to a standard normal reference. This follows R's
\code{summary.glm} convention. The \code{use\_t} option overrides it.
\code{wald\_chisq} is the squared statistic; \code{wald\_p\_value} uses its
asymptotic $\chi^2_1$ reference and can therefore differ from the displayed
$t$-based p-value. The \code{level} option defaults to 0.95.

\subsection{Availability and validation}
The inference object always provides \code{status} and \code{reason}.
Results may be \code{ok}, \code{disabled}, \code{skipped} by size, or
\code{unavailable}. Ordinary Wald results require an unpenalized fit whose final
gradient meets the configured tolerance, $n>p$, and nonsingular, sufficiently
well-conditioned statistical curvature. Detected separation, nonpositive
information, and degenerate dispersion prevent inference; the fitted model is
still returned. Separation screening checks the fitted predictor and is not an
exhaustive linear-programming separation test. Forcing inference bypasses only
the size policy. Ridge inference would require a separately defined approximation
or bias correction and is not reported as ordinary Wald inference.

For step-stopped fits, inspect convergence and, when appropriate, refit with
the following options in Python (use R/Julia boolean syntax):
\begin{lstlisting}
accept_stalled=False, accept_negligible=False
\end{lstlisting}
Original Julia exposes the stalled-step option;
the native Julia wrapper exposes both. No automatic refit is performed.

The retained benchmark scripts explicitly disable inference to keep fitting
and post-fit costs separate. Tests compare against R \code{glm}, statsmodels,
closed-form Gaussian results and independently computed sandwich matrices.
Relevant documentation: \href{https://stat.ethz.ch/R-manual/R-devel/library/stats/html/summary.glm.html}{R summary.glm},
\href{https://www.statsmodels.org/stable/generated/statsmodels.genmod.generalized_linear_model.GLMResults.html}{statsmodels GLMResults},
\href{https://juliastats.org/GLM.jl/latest/examples/}{GLM.jl examples}, and
\href{https://sandwich.r-forge.r-project.org/reference/sandwich.html}{sandwich covariance construction}.

\section{Reproducible Benchmarks}\label{sec:benchmark}
\subsection{Scope and data}
The retained comparison covers six GLM families---Gaussian identity,
Bernoulli logit, Bernoulli probit, Poisson log, Gamma log, and negative-binomial
log---and the four residual models. Gaussian-log, Gamma-inverse, Tweedie, and
Binomial are supported by the core but were not included in this timing study.

The archived timings predate the automatic-intercept interface update; they
are preserved measurements. Re-run the scripts to time the current code.
The archived study uses $n=2000$ and $p=1001$, including an intercept.
The Julia generator follows the original dense/sparse benchmark construction
and exports common binary matrices, responses, and initial coefficients.
Using the same binary inputs avoids cross-language differences in random-number
generation and decimal parsing. Rebuildable input files were removed during
cleanup; the generator and recorded results are retained.

\subsection{Run the comparison}
Julia, Python, and Rscript must be available on PATH. Install the Python/R
benchmark dependencies and build the native library first.
The commands below write to a fresh directory rather than overwriting the archive.
\begin{lstlisting}[language=bash]
julia --project=. -e 'using Pkg; Pkg.instantiate()'
BENCH_OUT=build/family-benchmark

julia --compiled-modules=existing --project=. \
  benchmark/native/family_performance.jl 1000 "$BENCH_OUT"

python3 benchmark/native/family_performance.py \
  --size 1000 --output "$BENCH_OUT" --language both

python3 benchmark/native/repeat_family_native.py \
  --output "$BENCH_OUT" --language both

python3 benchmark/native/report_family_performance.py \
  --output "$BENCH_OUT"
python3 benchmark/native/export_latex.py \
  --results "$BENCH_OUT"
\end{lstlisting}
The Julia step generates inputs and times the original Spectral-MM solver.
The first Python/R step includes applicable GLM references; the repeat step
measures the six native solvers in interleaved order.
Use \code{--language python} or \code{--language R} for one interface; the combined
report requires both interfaces' result files.

The argument 1000 is the number of features before adding the intercept, and the
sample size is twice that number. To run $n=10000$, $p=5001$, replace both size
arguments with 5000. Those larger dimensions are not the dimensions of the
archived results. GLM fits that reach the 1000-iteration limit can take a long time.

\subsection{Timing protocol and accuracy}
Each method is warmed up before timing. The initial comparison uses three
measured runs and reports their median. The native repeat uses five rounds with
shuffled method order; short calls are timed in batches. Thread settings are
restricted for the numerical libraries.

Timing includes fitting, wrapper conversion, and model construction. It excludes
imports, compilation, data generation, and scoring. The GLM references convert
sparse designs to dense arrays; conversion time is included.
Python uses \code{statsmodels.GLM}; R uses \code{glm.fit}, with a fixed
negative-binomial shape parameter supplied through MASS. No GLM reference is
substituted for the four residual models.

The benchmark uses $\rho=10^{-6}$, no ridge penalty, spectral rank 5, initial
Lanczos dimension 12, at most 1000 outer and 200 inner iterations,
\code{eta\_max}=0.6, \code{correction\_tol}=0.05,
\code{resid\_tol}=0.5, and outer Nesterov extrapolation.
The common gradient check uses \code{gtol}=$10^{-6}$ and
\code{relgtol}=$10^{-8}$. GLM implementations use their own stopping criteria
with an internal tolerance of $10^{-12}$ and are evaluated afterwards using the
common gradient measure.

Family parameters are fixed at $\theta=4$ for negative binomial,
$(\tau,\epsilon)=(0.25,0.1)$ for smooth quantile, $\tau=0.25$ for expectile,
$\delta=1$ for pseudo-Huber, and $(\nu,\sigma)=(4,1)$ for Student-$t$.

\subsection{Reading and exporting results}
The retained report is
\path{benchmark/native/results/family-performance/report.md}.
CSV files preserve the first-run and repeat summaries, individual timing samples,
and convergence diagnostics; configuration and environment records accompany them.
The main native tables use the interleaved repeat. GLM and Julia timings come
from the earlier three-run sessions, so cross-session timing differences should
not be interpreted as pure language overhead.

The binary examples exhibited separation/perfect prediction. Small gradients
therefore do not establish a unique finite maximum-likelihood estimate.
Several fits stopped on a step criterion without meeting the common gradient
threshold, and the Gamma GLM references reached the iteration limit.
These outcomes remain visible in the report rather than being counted as
same-accuracy speedups.

To compile the exported tables locally:
\begin{lstlisting}[language=bash]
cd "$BENCH_OUT/latex"
pdflatex -interaction=nonstopmode -halt-on-error performance_tables.tex
\end{lstlisting}
The table fragments require \code{booktabs}. A dagger marks solver-declared
convergence without passing the strict gradient check; a double dagger identifies
separable binary examples. Nonconverged runs retain their measured times and
reported diagnostics.

\section{Validation and Current Scope}
Completed validation compared all 14 families across native and Julia implementations,
including dense/sparse designs and all six solvers. Development test scripts and
generated fixtures were removed during cleanup; recorded validation results and
final benchmark data are retained.

The statistical interfaces provide matrix and formula fitting, standard result accessors,
and separate numerical diagnostics. Offsets and observation weights remain explicitly
unsupported by the numerical core. Optional inference is described above.
Native and Julia optimization trajectories need not be identical: Lanczos random
starts and floating-point operation orders differ. Agreement should be assessed
using the objective, gradient diagnostics, and fitted coefficients at a suitable
accuracy target.

\section*{Using This Document on Overleaf}
Upload this \code{.tex} file to a blank Overleaf project, set it as the main
document, and select \textbf{pdfLaTeX}. All required packages are standard LaTeX
packages. No external figures, bibliography database, shell escape, local fonts,
or additional input files are required. The commands in the code listings are
instructions to run in the local project, not commands executed by Overleaf.

\end{document}
